\documentclass[11pt]{article}
\usepackage[T1]{fontenc}
\usepackage[utf8]{inputenc}
\usepackage[a4paper,margin=1in]{geometry}
\usepackage{amsmath,amssymb,amsthm,mathtools}
\usepackage{xcolor}
\usepackage{hyperref}
\usepackage{enumitem}
\usepackage{setspace}
\usepackage{booktabs}
\usepackage{array}
\usepackage{graphicx}

\hypersetup{colorlinks=true,linkcolor=blue!50!black,urlcolor=blue!50!black,citecolor=blue!50!black}
\setlength{\parskip}{0.5em}
\setlength{\parindent}{0pt}
\onehalfspacing

% -- Theorem environments 
\newtheorem{theorem}{Theorem}[section]
\newtheorem{proposition}[theorem]{Proposition}
\newtheorem{lemma}[theorem]{Lemma}
\newtheorem{corollary}[theorem]{Corollary}
\newtheorem{definition}[theorem]{Definition}
\newtheorem{assumption}[theorem]{Assumption}
\newtheorem{observation}[theorem]{Observation}
\newtheorem*{remark}{Remark}
\newtheorem*{openproblem}{Open Problem}

% Custom label-free environments
\newenvironment{model}[1]{\par\smallskip\noindent\textit{[Model: #1]}\quad}{\par\smallskip}
\newenvironment{heuristic}[1]{\par\smallskip\noindent\textit{[Heuristic: #1]}\quad}{\par\smallskip}

% -- Macros 
\newcommand{\Hstr}{H_{\mathrm{str}}}
\newcommand{\Hnull}{H_{\mathrm{null}}}
\newcommand{\Estr}{E_{\mathrm{str}}}
\newcommand{\Espec}{E_{\mathrm{spec}}}
\newcommand{\etaorig}{\eta_{\mathrm{orig}}}
\newcommand{\Tren}{T_{\mathrm{ren}}}
\newcommand{\Gun}{G^{\mathrm{un}}}
\newcommand{\Gnorm}{G^{\mathrm{norm}}}
\newcommand{\acan}{a_p^{\mathrm{can}}}
\newcommand{\Dcan}{D^{\mathrm{can}}}
\newcommand{\Estrcan}{\Estr^{\mathrm{can}}}
\newcommand{\Especcan}{\Espec^{\mathrm{can}}}
\newcommand{\Hlocal}{H_{\mathrm{local}}}
\newcommand{\Hren}{H_{\mathrm{ren}}}
\newcommand{\cren}{c_p^{\mathrm{ren}}}

% -- Title 
\title{\vspace{-1.5em}\bfseries A Finite-Cutoff Hilbert-Space Model\\[0.4em]
\large for Prime--Zero Energy Structure}
\author{Ulrich Tehrani}
\date{Research Note \textperiodcentered\ March 2026 \textperiodcentered\ v1}

\begin{document}
\maketitle

\begin{abstract}
We construct a finite-cutoff Hilbert-space model for the spectral structure
associated with primes and nontrivial zeros of the Riemann zeta function.
Given a prime cutoff $\kappa$ and truncated zero ordinates
$\{\gamma_1,\ldots,\gamma_N\}$, we define two finite-dimensional real
Hilbert spaces $\Hstr$ (over primes) and $\Hnull$ (over zeros),
connected by a linear map $\Phi\colon\Hstr\to\Hnull$, and the
self-adjoint loop operator $T = \Phi^*\circ\Phi$.
We establish the algebraic identity
\[
\Delta := \Estr - \Espec
 = \sum_p c_p^2\|a_p\|^2 - \Bigl\|\sum_p c_p a_p\Bigr\|^2
\]
by direct expansion; numerically, $\Delta > 0$ and
$\etaorig = \Delta/\Estr \in [0.598,\,0.704]$ for benchmark parameters
$(\kappa,\varepsilon,N)=(53,0.05,100)$, $\sigma\in[0.1,0.95]$.
The formal core is self-contained in \S\S\,2--5;
results from Papers~1--2 appear only as context.
An analytic proof of $\etaorig>0$ remains open.
\emph{[Geometric interpretation, not used in any proof]:}
The model suggests a picture in which $\sigma=\tfrac12$ plays
the role of an energy reference point.
\end{abstract}

\tableofcontents
\newpage

%% ═══════════════════════════════════════════════════════════════════════════
\section{Introduction}
%% ═══════════════════════════════════════════════════════════════════════════

\subsection{Background and Motivation}

The present paper is the third in a series studying the curvature of the
Riemann zeta function from a geometric and operator-theoretic perspective.
A reader may read \S\S\,2--5 independently of the series context;
\S\S\,1.1--1.2 and \S\S\,7--10 provide motivation and interpretation
but are not needed for the formal results.
The paper is intended to be mathematically self-contained and readable
without prior knowledge of the earlier papers in the series.
Definitions and formal setup are in~\S2.
The proofs of the main formal statements are in \S\S\,3--5, and are
self-contained. Section~6 records open problems.
No external project material is needed for \S\S\,2--5.
The geometric interpretations in \S\S\,7--10 are clearly separated from
the formal content.

\textbf{Paper~1}~\cite{Paper1} introduced the function
\[
H_{\xi}(\sigma,t) := \frac{d^2}{d\sigma^2}\log|\xi(\sigma+it)|^2
\]
as a second-derivative decomposition of the log-completed zeta function, and
established its local positivity structure. A key finding was the divergence
$\Hlocal(\tfrac12,\kappa)\sim 8(\log\kappa)^2\to\infty$ at $\sigma=\tfrac12$,
contrasted with convergence for $\sigma>\tfrac12$. The connection to the Li
coefficients $\lambda_n$, whose positivity is equivalent to RH~\cite{Li97},
was identified.

\textbf{Paper~2}~\cite{Paper2} built on this by constructing explicit
Schwartz-class test functions $g^*\in\mathcal{S}_{\mathrm{ad}}$ adapted
to the adelic zeta integral. These functions realize the Weil explicit formula
as an identity
$W(g^*\ast\check{g}^*)=Z(g^*)-\Hlocal(\sigma,\kappa)+O(\varepsilon)$,
providing a bridge between local curvature data and the Weil positivity
criterion.

\textbf{The present paper} makes the geometric structure underlying both
papers explicit. The central objects --- two finite-dimensional Hilbert spaces
connected by a linear map $\Phi$, and a self-adjoint loop operator ---
are motivated by structures appearing in Papers~1 and~2 and are here
constructed and analyzed directly.

\subsection{Imported Results from Papers 1 and 2}
\label{sec:imported}

This paper builds on the following results established in Papers~1 and~2:

\begin{itemize}[leftmargin=2.2em]
\item[\textbf{(R1)}] The truncated local curvature
$\Hlocal(\sigma,\kappa)=\psi_1(\sigma)+\sum_{p\le\kappa}V_p(\sigma)$ satisfies
$\Hlocal(\tfrac12,\kappa)\sim 8(\log\kappa)^2\to\infty$, while
$\Hlocal(\sigma,\kappa)\to C(\sigma)<\infty$ for all $\sigma>\tfrac12$.
\emph{[Paper~1, Lemma~2]}

\item[\textbf{(R2)}] An explicit family of admissible test functions $g^*$
satisfies
$W(g^*\ast\check{g}^*)=Z(g^*)-\Hlocal(\sigma,\kappa)+O(\varepsilon)$.
\emph{[Paper~2, Theorem~3.1]}

\item[\textbf{(R3)}] The renormalized prime weights $\cren=f_p^{1/2}>0$ satisfy
$D=\sum_p(\cren)^2\approx 9.471<\infty$. \emph{[Paper~2, \S4]}
\end{itemize}

\textbf{Reader contract:} Results R1--R3 are not used in any proof or
theorem-level argument in \S\S\,2--5. They appear in \S\,2.3 (Remark~2.3)
for motivation only, and formally in \S\,10.2 for contextual discussion.
\S\S\,2--5 are a finite-dimensional operator model; they do not require
acceptance of any RH-related interpretation.

\subsection{What This Paper Proves and What It Imports}

\begin{center}
\begin{tabular}{lll}
\toprule
Content & Status & Where \\
\midrule
$\Hstr, \Hnull, \Phi, T$ (\S\S\,2--5) & Proved here & \S\S\,2--5 \\
Algebraic identity $\Delta=\Estr-\Espec$ & Proved here & \S4 \\
$\etaorig\in[0.598,0.704]$ & Numerical (this paper) & \S5 \\
$\Hlocal(\tfrac12,\kappa)\sim 8(\log\kappa)^2$ (R1) & Imported \cite{Paper1} & \S\,2.3, \S\S\,8--9 \\
$W(\cdot)=Z(g^*)-\Hlocal+O(\varepsilon)$ (R2) & Imported \cite{Paper2} & \S\,10.2 \\
$\cren=f_p^{1/2}$, $D\approx 9.471$ (R3) & Imported \cite{Paper2} & \S\,10.2 \\
\S\S\,7--10 geometric model & Interpretive only & Not in proofs \\
\bottomrule
\end{tabular}
\end{center}

\subsection{Connection to Papers 1 and 2}

The present paper isolates a finite-cutoff Hilbert-space structure underlying
the energy comparisons of Papers~1 and~2, while the full normalization-level
identification with Paper~2 remains open (\S\,10.2). The renormalized
framework ($Z_{\mathrm{ren}}, \Hren, \cren$) of Paper~2 and the
$\eta$-framework ($\Estr, \Espec, c_p$) of this paper use different
normalizations; a precise normalization-level identification remains open
(see \S\,10.2 and Open Problem~6.5).

\subsection{Relation to Existing Work}

The construction of $T=\Phi^*\circ\Phi$ from prime resonance vectors provides
a concrete Hilbert-space realization outside the Hilbert-P\'olya setting:
$T$ is explicitly constructed from prime-resonance vectors,
and its eigenvalues are not the zero ordinates.
The adelic viewpoint connects to Connes' noncommutative geometry
approach~\cite{Con99}, and the test functions of Paper~2 have the structure
of adelic Schwartz functions as in Tate's thesis~\cite{Tat50}. The Weil
positivity criterion~\cite{Wei52,Lag99} provides the outer framework.

\subsection{Outline}

Section~2 defines the two Hilbert spaces and their inner products.
Section~3 constructs the linear map $\Phi$ and the loop operator $T$.
Section~4 proves the variance identity and defines $\etaorig$.
Section~5 collects numerical results.
Section~6 states open problems.
Sections~7--10 develop the geometric interpretation (all statements in these
sections are explicitly marked as geometric model or heuristic).

%% ═══════════════════════════════════════════════════════════════════════════
\section{Two Hilbert Spaces}
%% ═══════════════════════════════════════════════════════════════════════════

We work throughout with real Hilbert spaces; all vectors and operators have
real entries. Complex phasor language appears in \S\,3.1 and \S\S\,7--10 as
heuristic motivation only; all formal objects ($\Hstr, \Hnull, \Phi, T$)
are real.

Throughout this paper, $\kappa>1$ is a fixed real cutoff parameter,
$S(\kappa):=\{p\text{ prime}:p\leq\kappa\}$ the set of active primes,
and $\{\gamma_k\}_{k=1}^N$ the first $N$ positive imaginary parts of
nontrivial zeros of $\zeta(s)$, ordered $0<\gamma_1\leq\gamma_2\leq\cdots$.
The Gaussian damping parameter $\varepsilon>0$ is fixed throughout.

\begin{definition}[Prime Space $\Hstr$]
\[
\Hstr := \ell^2(S(\kappa))
\]
with basis $\{e_p : p\in S(\kappa)\}$, standard inner product
$\langle a,b\rangle_{\mathrm{str}}=\sum_p a_p b_p$, and
$\dim\Hstr=|S(\kappa)|=\pi(\kappa)$.
The \emph{structural energy} of a weight vector $c\in\Hstr$ is
$\Estr(c)=\sum_p c_p^2\|a_p\|^2$ (defined in \S4 after $\Phi$ is introduced).
\end{definition}

\begin{definition}[Zero Space $\Hnull$]
\[
\Hnull := \ell^2(\{\gamma_1,\ldots,\gamma_N\})
\]
with basis $\{f_k:k=1,\ldots,N\}$, inner product
$\langle u,v\rangle_{\mathrm{null}}=\sum_{k=1}^N u_k v_k$
(unweighted), and $\dim\Hnull=N$.
\end{definition}

\begin{remark}[Why $\sigma=\tfrac12$ is a distinguished reference point]
\label{rem:half}
\emph{[Context only --- not used in any argument in \S\S\,3--5.]}

Two independent proven reasons support this choice:

\emph{(i) Symmetry} [proven]: $\sigma=\tfrac12$ is the unique fixed point of
$s\mapsto 1-s$, the symmetry of the completed zeta function
$\xi(s)=\xi(1-s)$.

\emph{(ii) Curvature singularity} [proven, Papers~1--2]:
$\Hlocal(\tfrac12,\kappa)\sim 8(\log\kappa)^2\to\infty$ as $\kappa\to\infty$,
while $\Hlocal(\sigma,\kappa)\to C(\sigma)<\infty$ for all $\sigma>\tfrac12$.

$\sigma=\tfrac12$ is thus a distinguished reference point by symmetry and
curvature divergence. Energy near-equality ($\Estr\approx\Espec$) is NOT
observed in the finite-cutoff model (see \S5 and \S8).
\end{remark}

%% ═══════════════════════════════════════════════════════════════════════════
\section{The Linear Map $\Phi$ and the Loop Operator $T$}
%% ═══════════════════════════════════════════════════════════════════════════

\textbf{Convention (Option B throughout):} Damping lives in the vectors:
$a_p=(e^{-\varepsilon^2\gamma_k^2/2}\sin(\gamma_k\log p))_k$.
The $\Hnull$ inner product is unweighted: $\langle u,v\rangle=\sum_k u_kv_k$.

\subsection{Motivation for the resonance vectors}

\emph{[Heuristic motivation --- the formal definition is Definition~3.2 below.]}

The prime resonance vector at zero ordinate $\gamma_k$ and prime $p$ arises as
the imaginary part of the prime phasor $\varphi_p(u)=e^{iu\log p}$ evaluated
at $u=\gamma_k$, with Gaussian regularization for convergence:
\[
(a_p)_k = e^{-\varepsilon^2\gamma_k^2/2}\cdot\sin(\gamma_k\log p).
\]
The $k$-th component measures how strongly prime $p$ resonates with zero
ordinate $\gamma_k$.

\begin{definition}[Prime Resonance Vectors]
For each prime $p\in S(\kappa)$, the \emph{canonical resonance vector}
at $\sigma=\tfrac12$ is
\[
\acan := a_p :=
\bigl(e^{-\varepsilon^2\gamma_k^2/2}\sin(\gamma_k\log p)\bigr)_{k=1}^N
\;\in\;\Hnull.
\]
At general $\sigma$, the resonance vectors depend on $\sigma$ via
\[
a_p^{(\sigma)} := p^{-(\sigma-1/2)}\cdot\acan.
\]
The vectors $a_p=\acan=a_p^{(1/2)}$ at $\sigma=\tfrac12$ are the
canonical representatives used in all primary definitions.
\end{definition}

\begin{definition}[Linear Map $\Phi$]
\[
\Phi\colon\Hstr\to\Hnull,\qquad\Phi(e_p):=\acan.
\]
Extended linearly: $\Phi(c)=\sum_p c_p\acan$ for any
$c=\sum_p c_p e_p\in\Hstr$.

\emph{We define $\Phi$ as a linear map; no injectivity claim is made for the
finite-cutoff model.}
\end{definition}

\begin{definition}[Gram Matrices]
The \emph{unnormalized Gram matrix} (canonical) is
\[
\Gun_{pq} := \Gun{}^{\mathrm{,can}}_{pq}
:= \langle\acan,a_q^{\mathrm{can}}\rangle_{\Hnull}
= \sum_{k=1}^N e^{-\varepsilon^2\gamma_k^2}\sin(\gamma_k\log p)\sin(\gamma_k\log q).
\]
The $\sigma$-dependent extension is
$\Gun(\sigma)_{pq}=p^{-(\sigma-1/2)}q^{-(\sigma-1/2)}\Gun_{pq}$.
The \emph{normalized Gram matrix} is
\[
\Gnorm_{pq}:=\frac{\Gun_{pq}}{\|\acan\|\cdot\|a_q^{\mathrm{can}}\|}
=\Bigl\langle\frac{\acan}{\|\acan\|},\frac{a_q^{\mathrm{can}}}{\|a_q^{\mathrm{can}}\|}\Bigr\rangle_{\Hnull}.
\]
The matrix $\Gnorm$ is $\sigma$-invariant (Proposition~3.5(iv)).
\end{definition}

\begin{definition}[Loop Operator $T$]
\[
T:=\Phi^*\circ\Phi\colon\Hstr\longrightarrow\Hstr.
\]
The matrix representation of $T$ in the prime basis is $T_{pq}=\Gun_{pq}$.
$T$ is thus identical to the canonical Gram matrix $\Gun$, and is explicitly
constructed from prime resonance vectors.
\end{definition}

\begin{proposition}[Properties of $T$]
\label{prop:T}
\begin{enumerate}[label=\emph{(\roman*)}, leftmargin=2.2em]
\item \emph{Self-adjoint:} $T^*=(\Phi^*\Phi)^*=\Phi^*\Phi=T$.
\item \emph{Positive semi-definite:} $\langle Tc,c\rangle=\|\Phi c\|^2\geq 0$
for all $c\in\Hstr$.
\item \emph{Non-negative eigenvalues:} All eigenvalues $\lambda_j\geq 0$.
\item \emph{$\sigma$-invariance of $\Gnorm$} [by direct calculation]:
For the $\sigma$-dependent family $a_p^{(\sigma)}$, the normalized Gram matrix
$\Gnorm_{pq}$ is independent of $\sigma$. \emph{Proof:} The scalar factor
$p^{-(\sigma-1/2)}$ cancels in the normalized ratio. \qed
\end{enumerate}
\end{proposition}

\begin{remark}[Spectral structure of $T$]
The eigenvalues $\lambda_j\geq 0$ of $T=\Phi^*\circ\Phi$ are the squares of
the singular values $\sigma_j(\Phi)$: $\lambda_j=\sigma_j(\Phi)^2$.
The eigenvalues $\lambda_j$ are built from the zero ordinates $\gamma_k$
appearing in the definition of $\acan$, but are \emph{not} directly the zeros
$\gamma_k$ themselves. This distinguishes $T$ from the classical Hilbert-P\'olya
ansatz: $T$ is explicitly constructed, not postulated.
\end{remark}

%% ═══════════════════════════════════════════════════════════════════════════
\section{The Variance Identity and $\etaorig$}
%% ═══════════════════════════════════════════════════════════════════════════

\textbf{Convention:} Objects without $\sigma$-argument refer to the canonical
($\sigma=\tfrac12$) version throughout \S\S\,4--5.

\begin{definition}[Canonical and $\sigma$-dependent Energies]
\label{def:energies}
For a weight vector $c\in\Hstr$ with components $c_p$:

\emph{Canonical quantities} (at $\sigma=\tfrac12$):
\[
\Dcan:=\mathrm{diag}\!\bigl(\|\acan\|^2\bigr)_{p\in S(\kappa)},
\qquad
\Estrcan:=c^\top\Dcan c=\sum_p c_p^2\|\acan\|^2,
\]
\[
\Especcan:=\|\Phi c\|^2=c^\top\Gun c.
\]

\emph{$\sigma$-dependent extensions:}
\[
D(\sigma)_{pp}:=p^{-2(\sigma-1/2)}\cdot\Dcan_{pp},
\qquad
\Estr(\sigma):=\sum_p c_p^2\,p^{-2(\sigma-1/2)}\|\acan\|^2,
\]
\[
\Espec(\sigma):=c^\top\Gun(\sigma)\,c.
\]
The \emph{canonical weight vector} in the $\eta$-context is
\[
c_p=\sqrt{8}\cdot\frac{\log p}{p},\qquad p\in S(\kappa).
\]
\end{definition}

\begin{assumption}[Positive definiteness of $D$]
Throughout this paper we assume $D_{pp}=\|a_p\|^2>0$ for all
$p\in S(\kappa)$. This holds in all numerical experiments conducted here.
\end{assumption}

\begin{theorem}[Algebraic Identity]
\label{thm:identity}
The following equality holds by direct expansion:
\[
\Delta:=\Estr-\Espec=\sum_p c_p^2\|a_p\|^2-\Bigl\|\sum_p c_p a_p\Bigr\|^2.
\]
\end{theorem}

\begin{proof}
Expand $\Espec=\|\Phi c\|^2=\bigl\langle\sum_p c_p a_p,\sum_q c_q
a_q\bigr\rangle=\sum_{p,q}c_p c_q\langle a_p,a_q\rangle=c^\top\Gun c$.
Then $\Delta=c^\top\Dcan c-c^\top\Gun c$, which is the stated identity.
\end{proof}

\begin{observation}[Non-negativity, numerical]
\label{obs:nonneg}
\[
\Delta\geq 0\quad\text{equivalently}\quad\etaorig\geq 0
\]
holds for all tested parameters with canonical weights
$c_p=\sqrt{8}\log(p)/p$ (Definition~\ref{def:energies}).
\emph{[Numerical only --- no unconditional proof for general $c$.]}

Note: $\etaorig\geq 0$ for general weight vector $c$ is equivalent to
$\lambda_{\max}(D^{-1/2}TD^{-1/2})\leq 1$, which is the content of
Open Problem~6.2. Numerical experiments show
$\lambda_{\max}(D^{-1/2}TD^{-1/2})\approx 0.39\cdot\pi(\kappa)\to\infty$,
so the universal bound $\lambda_{\max}\leq 1$ does not hold in general.

\emph{Cauchy-Schwarz does \textbf{not} guarantee $\etaorig\geq 0$}: the
correct identity is $\Delta=c^\top(\Dcan-\Gun)c$; the sign of $\Delta$ is not
determined by the algebraic identity alone.
\end{observation}

\begin{definition}[$\etaorig$]
\label{def:eta}
\[
\etaorig:=1-\frac{\Espec}{\Estr}=1-\frac{\|\Phi c\|^2}{\Estr}.
\]
Boundary cases: $\etaorig=0$ iff $\Espec=\Estr$; $\etaorig=1$ iff $\Phi c=0$.
\end{definition}

\begin{remark}[$\sigma$-dependence via rescaled weights]
The entire $\sigma$-dependence of $\etaorig(\sigma)$ is carried by the
rescaled weight vector $\tilde{c}_p(\sigma):=c_p\cdot p^{-(\sigma-1/2)}$:
\[
\etaorig(\sigma)=1-\frac{\tilde{c}(\sigma)^\top\Gun\,\tilde{c}(\sigma)}
{\tilde{c}(\sigma)^\top\Dcan\,\tilde{c}(\sigma)}.
\]
\emph{Proof:} Since $\Gun(\sigma)_{pq}=p^{-(\sigma-1/2)}q^{-(\sigma-1/2)}\Gun_{pq}$,
one has $c^\top\Gun(\sigma)c=\tilde{c}^\top\Gun\tilde{c}$ and
$\Estr(\sigma)=\tilde{c}^\top\Dcan\tilde{c}$. \qed
\end{remark}

%% ═══════════════════════════════════════════════════════════════════════════
\section{Numerical Results}
%% ═══════════════════════════════════════════════════════════════════════════

\textbf{Setup.}
Zeros $\gamma_k$ computed via \texttt{mpmath.zetazero()}, $N=100$ throughout.
$\sigma$~grid: $\{0.1,0.2,\ldots,0.9,0.95\}$.
Eigenvectors and eigenvalues: NumPy \texttt{linalg.eigh()}.
Requirements: mpmath $\geq 1.3$; NumPy $\geq 1.24$.
All computations via \texttt{eta\_verification.py}; companion code at
\url{https://github.com/utehrani/analysislab-nt}.
Algebraic identity (Theorem~\ref{thm:identity}) numerically verified to
rounding error $<10^{-14}$.

\emph{All results in this section are numerical. Parameters tested:
$\kappa\in\{23,53,101,199\}$, $\varepsilon\in\{0.02,0.05,0.10\}$,
$\sigma\in[0.1,0.95]$, $N=100$.}

\begin{observation}[$\etaorig$ positive across all tested parameters]
$\etaorig>0$ in all tested parameter combinations. For the benchmark slice
$(\kappa,\varepsilon,N)=(53,0.05,100)$, $\sigma\in[0.1,0.95]$:
\[
\etaorig\in[0.598,\;0.704].
\]
The function $\sigma\mapsto\etaorig(\sigma)$ is strictly monotone decreasing.
Note: the range $[0.598,0.704]$ is strongly $\varepsilon$-dependent; no
universal range over all $(\kappa,\varepsilon)$ is claimed.
\end{observation}

\textbf{Table 5.2}: $\etaorig(\sigma)$ for $(\kappa,\varepsilon,N)=(53,0.05,100)$.
[Numerical]

\begin{center}
\begin{tabular}{cc}
\toprule
$\sigma$ & $\etaorig$ \\
\midrule
0.10 & 0.70441977 \\
0.20 & 0.69826865 \\
0.30 & 0.69038398 \\
0.40 & 0.68072993 \\
0.50 & 0.66926873 \\
0.60 & 0.65602069 \\
0.70 & 0.64110102 \\
0.80 & 0.62472010 \\
0.90 & 0.60715692 \\
0.95 & 0.59802896 \\
\bottomrule
\end{tabular}
\end{center}

\begin{observation}[Uniform spectral bound $B_{\max}\leq 0.106<1$]
Define the mode contribution
$B_j(\sigma):=\lambda_j(\sigma)|\langle c,u_j(\sigma)\rangle|^2/\Estr(\sigma)$
with consistency identity $\sum_j B_j(\sigma)=1-\etaorig(\sigma)$ [algebraic].
The empirical diagnostic $B_{\max}(\sigma):=\max_{j:\lambda_j>1}B_j(\sigma)$
satisfies
\[
B_{\max}(\kappa,\varepsilon,\sigma)\leq 0.106<1
\]
for all tested $(\kappa,\varepsilon,\sigma)$. Selected values:

\begin{center}
\begin{tabular}{cccc}
\toprule
$\kappa$ & $\varepsilon=0.02$ & $\varepsilon=0.05$ & $\varepsilon=0.10$ \\
\midrule
23  & 0.1047 & 0.0220 & 0.0000 \\
53  & 0.0793 & 0.0737 & 0.0000 \\
101 & 0.0537 & 0.0844 & 0.0324 \\
199 & 0.0891 & 0.0461 & 0.1058 \\
\bottomrule
\end{tabular}
\end{center}
\end{observation}

\begin{observation}[$B_{\max}(\sigma)$ numerically $\sigma$-invariant]
For $(\kappa,\varepsilon)=(53,0.05)$, $B_{\max}(\sigma)=0.0737$ for all
tested $\sigma\in\{0.3,0.4,0.5,0.6,0.7\}$. Analytical explanation: open
(Open Problem~6.4).
\end{observation}

\begin{remark}[Cancellation mechanism, numerical analysis]
The mechanism underlying $B_{\max}<1$ is a cancellation between the monotone
weight vector $c_p=\sqrt{8}\log(p)/p$ and the oscillatory eigenvectors $u_j$
of $\Gun$: for $\kappa=53$, eigenvectors $u_1,\ldots,u_4$ each exhibit 10
sign changes along the prime ordering. Cancellation rates 67--97\% observed
for $\kappa\in\{23,53,101,199\}$. Formal proof: Open Problem~6.3.
\end{remark}

%% ═══════════════════════════════════════════════════════════════════════════
\section{Open Problems}
%% ═══════════════════════════════════════════════════════════════════════════

\textbf{(a) Within the $\eta$-framework of this paper:}

\begin{openproblem}[Analytic positivity]
Prove analytically that $\etaorig>0$ for the canonical weight vector
$c_p=\sqrt{8}\log(p)/p$.
Numerically confirmed: $\etaorig\in[0.598,0.704]$ for $(\kappa,\varepsilon,N)=(53,0.05,100)$,
$\sigma\in[0.1,0.95]$ (Table~5.2).
\end{openproblem}

\begin{openproblem}[Spectral bound for canonical weights]
Let $\Tren:=D^{-1/2}TD^{-1/2}$.
Numerical experiments show $\lambda_{\max}(\Tren)\approx 0.39\cdot\pi(\kappa)\to\infty$;
the universal bound $\lambda_{\max}<1$ (for all $c$) does not hold.
However, for the canonical weight vector $c_p=\sqrt{8}\log(p)/p$,
the quantity $\langle\Tren\tilde{c},\tilde{c}\rangle$ satisfies
$\langle\Tren\tilde{c},\tilde{c}\rangle<1$ (equivalently $\etaorig>0$)
for all tested $\kappa\in\{23,53,101,199,503,1009\}$.
Prove this analytically. The canonical $c_p$ appears orthogonal to the
dominant eigenmode of $\Tren$ (localized at $\gamma_1=14.13\ldots$);
proving this orthogonality is the core open problem.
\end{openproblem}

\begin{openproblem}[Analytical cancellation bound]
Prove analytically:
\[
\max_{j:\,\lambda_j>1}\frac{|\langle c,u_j\rangle|^2}{\|c\|^2}\leq C_1\ll 1
\quad\text{uniformly in }\kappa,j.
\]
Approach: Abel summation for $c_p\sim(\log p)/p$ against oscillatory partial
sums $\sum_{q\leq p}(u_j)_q$, without reference to RH.
Numerically: $C_1<0.36$ for $(\kappa,\varepsilon)=(53,0.05)$.
\end{openproblem}

\begin{openproblem}[$\sigma$-invariance of $B_{\max}$ analytically]
The numerical $\sigma$-invariance of $B_{\max}(\sigma)$ (consistent with
Proposition~\ref{prop:T}(iv)) is not yet derived analytically from the weight
rescaling $c\mapsto\tilde{c}(\sigma)$.
\end{openproblem}

\textbf{(b) Beyond the scope of this paper:}

\begin{openproblem}[Renormalized shell energy]
The renormalized shell energy $\Hren(\kappa,\varepsilon(\kappa))$ for canonical
coupling $N_{\mathrm{eff}}=\pi(\kappa)$ is not yet computed. This lies in the
Paper~2 normalization framework (see \S\,10.2).
\end{openproblem}

\begin{remark}[Connection to the Riemann Hypothesis]
\emph{[Geometric interpretation --- not a theorem.]}
In the geometric model, RH is modeled heuristically as injectivity of an
idealized linear map $\Phi$. No formal connection between injectivity of the
finite-cutoff $\Phi$ and $\etaorig>0$ is established here. The quantitative
convexity inequality of Papers~1--2 is an open problem in the earlier framework
and is \textbf{not} addressed in this paper.
\end{remark}

%% ═══════════════════════════════════════════════════════════════════════════
\section{Four Interpretive Levels (Geometric Interpretation)}
%% ═══════════════════════════════════════════════════════════════════════════

\textbf{Notation for \S\S\,7--10.}
Every statement using model labels
(\emph{prime resonance, interference field, resonance frequencies, energy
reference point}) is a geometric interpretation --- not a theorem.
Formal objects retain their definitions from \S\S\,2--5.
No map ($\Phi_{01},\Phi_{12},\Phi_{23}$), no intermediate space $\mathcal{I}$,
and no duality statement in \S\S\,7--10 is part of the theorem-level content.

\textbf{[Model 7.1]} (Four Interpretive Levels).
We organize the arithmetic and spectral structures into four levels:
\begin{align*}
\Omega_0 &: \mathbb{N}\quad\text{(arithmetic, discrete)}\\
\Omega_1 &: \Hstr\quad\text{(prime exponent lattice, structural energy)}\\
\Omega_2 &: \mathcal{I}\quad\text{(interference fields)}\\
\Omega_3 &: \Hnull\quad\text{(spectral space, zero resonance frequencies)}
\end{align*}
with model maps (mnemonics only, not formal objects):
$\Phi_{01}:\mathbb{N}\to\Hstr$,
$\Phi_{12}:\Hstr\to\mathcal{I}$,
$\Phi_{23}:\mathcal{I}\to\Hnull$.

\textbf{[Heuristic 7.3]} (Zeros as Resonance Frequencies).
In this geometric picture, the zeros $\gamma_k$ may be viewed as resonance
frequencies of the prime interference field: the condition
$\zeta(\sigma+i\gamma_k)=0$ is modeled as destructive interference of the
prime resonance waves $p^{-\sigma-i\gamma_k}$ at frequency $\gamma_k$.
\emph{This is a heuristic model picture.}

%% ═══════════════════════════════════════════════════════════════════════════
\section{The Energy Plane}
%% ═══════════════════════════════════════════════════════════════════════════

\textbf{[Model 8.1]} (Energy Coordinates).
For canonical weight vector $c$ and parameter $\sigma$, define
$X(\sigma):=\sqrt{\Estr(\sigma)}$ and $Y(\sigma):=\sqrt{\Espec(\sigma)}$.

\textbf{[Numerically verified, $(\kappa,\varepsilon,\sigma)=(53,0.05,0.5)$]:}
Energy near-equality does NOT hold in the finite-cutoff model:
\[
\Espec/\Estr\approx 0.331\qquad(\text{not }\approx 1).
\]
$\sigma=\tfrac12$ remains distinguished by proven properties (symmetry and
curvature divergence), NOT by energy equilibrium.

\textbf{[Model 8.3]} (Variance Identity in the Energy Plane).
$\Delta=\Estr-\Espec=X^2-Y^2=(X+Y)(X-Y)$.
The condition $\etaorig>0$ is represented as $X>Y$, observed numerically
throughout (Observation~\ref{obs:nonneg}).

%% ═══════════════════════════════════════════════════════════════════════════
\section{Why $\sigma=\tfrac{1}{2}$ is Special: Four Reasons}
%% ═══════════════════════════════════════════════════════════════════════════

\emph{Reason~(i) is proven. Reason~(ii) is numerical. Reason~(iii) is proven.
Reason~(iv) is a conceptual model only.}

\textbf{(i) Symmetry} [proven]:
$\sigma=\tfrac12$ is the unique fixed point of $s\mapsto 1-s$,
the symmetry underlying $\xi(s)=\xi(1-s)$.

\textbf{(ii) Energy dominance} [numerical, \S5]:
$\Estr(\tfrac12)>\Espec(\tfrac12)$ for all tested $(\kappa,\varepsilon)$,
with $\Espec/\Estr\approx 0.331$ at $(\kappa,\varepsilon)=(53,0.05)$.
For the benchmark slice, $\etaorig(\sigma)$ is strictly monotone decreasing
on $[0.1,0.95]$ (Table~5.2). Energy near-equality
($\Estr\approx\Espec$) is NOT observed.

\textbf{(iii) Curvature singularity} [proven, Papers~1--2]:
\[
\Hlocal(\tfrac12,\kappa)\sim 8(\log\kappa)^2\to\infty,\qquad
\Hlocal(\sigma,\kappa)\to C(\sigma)<\infty\text{ for }\sigma>\tfrac12.
\]

\textbf{(iv)} [Conceptual model, NOT numerically verified]: In the energy-plane
model of \S8, one might expect $\sigma=\tfrac12$ to be the
``45${}^\circ$-point'' where $X(\sigma)\approx Y(\sigma)$ in an idealized
infinite model. This is not confirmed in the finite-cutoff numerics:
$\sqrt{\Espec/\Estr}\approx 0.575\neq 1$ at $(\kappa,\varepsilon)=(53,0.05)$.

%% ═══════════════════════════════════════════════════════════════════════════
\section{Duality $\mathbb{N}\leftrightarrow\mathbb{C}$ and Outlook}
%% ═══════════════════════════════════════════════════════════════════════════

\textbf{[Model 10.1]} (Structural Duality).
\begin{center}
\begin{tabular}{ll}
\toprule
\textbf{Structural side} $\Omega_1$ & \textbf{Spectral side} $\Omega_3$ \\
\midrule
Basis: primes $p$ & Basis: zero ordinates $\gamma_k$ \\
Energy: $\Estr$ & Energy: $\Espec$ \\
Structure: multiplication & Structure: interference \\
Tool: Euler product & Tool: Explicit formula \\
\bottomrule
\end{tabular}
\end{center}

\textbf{[Model 10.2]} (Connection to Weil Positivity).
\emph{[Contextual only --- not part of the formal proof core \S\S\,2--6.]}

In Paper~2, the Weil functional $W(g^*\ast\check{g}^*)$, whose positivity for
all admissible $g^*$ is equivalent to RH~\cite{Wei52,Lag99}, satisfies
$W(g^*\ast\check{g}^*)=Z(g^*)-\Hlocal(\sigma,\kappa)+O(\varepsilon)$.
This suggests an energy-balance analogy with the $\eta$-framework, but the
$\eta$-framework and Weil framework live in different normalizations.
No theorem in this paper identifies $(\Estr,\Espec,\Delta)$ with
$(Z_{\mathrm{ren}},\Hren,Z_{\mathrm{off}})$ (see Open Problem~6.5).

\textbf{Outlook 10.3.}
Direct next problems: (1) analytic proof of $\etaorig>0$ via uniform partial
summation (Open Problem~6.3); (2) spectral bound for canonical weights (Open
Problem~6.2); (3) renormalized shell energy $\Hren(\kappa,\varepsilon(\kappa))$
(Open Problem~6.5). Broader questions include the connection to Connes'
program~\cite{Con99} and whether the test function family $\{g^*_{\sigma,\varepsilon}\}$
is a complete witness system for the Weil criterion~\cite{Lag99}.

%% ═══════════════════════════════════════════════════════════════════════════
\section*{References}
%% ═══════════════════════════════════════════════════════════════════════════

\begin{thebibliography}{99}

\bibitem{Con99}
Connes, A.
\textit{Trace formula in noncommutative geometry and the zeros of the Riemann
zeta function.}
Selecta Math.\ (N.S.)\ \textbf{5} (1999), 29--106.

\bibitem{Edw74}
Edwards, H.M.
\textit{Riemann's Zeta Function.}
Academic Press, 1974.

\bibitem{IK04}
Iwaniec, H.; Kowalski, E.
\textit{Analytic Number Theory.}
AMS, 2004.

\bibitem{Lag99}
Lagarias, J.C.
\textit{On a positivity property of the Riemann $\xi$-function.}
Acta Arith.\ \textbf{89} (1999), 217--234.

\bibitem{Li97}
Li, X.-J.
\textit{The positivity of a sequence of numbers and the Riemann hypothesis.}
J.\ Number Theory \textbf{65} (1997), 325--333.

\bibitem{Odl92}
Odlyzko, A.M.
\textit{The $10^{20}$-th zero of the Riemann zeta function and 70 million of
its neighbors.} Unpublished, 1992.
\url{https://www-users.cse.umn.edu/~odlyzko/unpublished/zeta.10to20.1992.pdf}

\bibitem{Paper1}
Tehrani, U.
\textit{A Curvature Decomposition of the Explicit Formula for the Riemann Zeta
Function.}
Zenodo, \href{https://doi.org/10.5281/zenodo.19025598}{doi:10.5281/zenodo.19025598},
v7, March 2026.

\bibitem{Paper2}
Tehrani, U.
\textit{From Local Curvature to the Weil Functional: An Explicit Construction.}
Zenodo, \href{https://doi.org/10.5281/zenodo.19106992}{doi:10.5281/zenodo.19106992},
v4, March 2026.

\bibitem{Sar05}
Sarnak, P.
\textit{Problems of the Millennium: The Riemann Hypothesis.}
Clay Math.\ Institute, 2004.
In: \textit{The Millennium Prize Problems}, AMS, 2006.

\bibitem{Tat50}
Tate, J.T.
\textit{Fourier analysis in number fields, and Hecke's zeta-functions.}
In: \textit{Algebraic Number Theory} (Cassels--Fr\"ohlich), Academic Press, 1967.

\bibitem{Wei52}
Weil, A.
\textit{Sur les ``formules explicites'' de la th\'eorie des nombres premiers.}
Comm.\ S\'em.\ Math.\ Univ.\ Lund (1952), 252--265.

\end{thebibliography}

%% ═══════════════════════════════════════════════════════════════════════════
\appendix
%% ═══════════════════════════════════════════════════════════════════════════

\section{Variable Reference}

\begin{center}
\renewcommand{\arraystretch}{1.25}
\begin{tabular}{p{3.4cm}p{7.0cm}p{2.0cm}}
\toprule
Symbol & Definition & Status \\
\midrule
$\kappa$ & prime cutoff & parameter \\
$S(\kappa)$ & primes $p\leq\kappa$ & definition \\
$\gamma_k$ & $k$-th zero ordinate (positive imag.\ part) & input \\
$\varepsilon$ & Gaussian damping & parameter \\
$\acan=a_p$ & canonical resonance vector at $\sigma=\tfrac12$ & \S3.1 \\
$a_p^{(\sigma)}$ & $=p^{-(\sigma-1/2)}\acan$ & \S3.1 \\
$\Gun_{pq}$ & canonical unnormalized Gram matrix & \S3.3 \\
$\Gun(\sigma)_{pq}$ & $\sigma$-dependent Gram & \S3.3 \\
$\Gnorm_{pq}$ & normalized Gram; $\sigma$-invariant [proven] & \S3.3 \\
$T=\Phi^*\Phi$ & loop operator $=\Gun$ & \S3.4 \\
$\Dcan$ & $\mathrm{diag}(\|\acan\|^2)$ & \S4.1 \\
$\Estrcan$ & $=c^\top\Dcan c$ & \S4.1 \\
$\Especcan$ & $=\|\Phi c\|^2=c^\top\Gun c$ & \S4.1 \\
$\Delta$ & $=\Estr-\Espec$ [algebraic]; sign numerical & \S4.2a/b \\
$\etaorig$ & $=\Delta/\Estr$ & \S4.3 \\
$\tilde{c}_p(\sigma)$ & $=c_p\cdot p^{-(\sigma-1/2)}$ & \S4.4 \\
$c_p$ & $=\sqrt{8}\log(p)/p$ (canonical weight) & \S4.1 \\
$B_j(\sigma)$ & $=\lambda_j|\langle c,u_j\rangle|^2/\Estr$; $\sum_jB_j=1-\etaorig$ & \S5.1 \\
$B_{\max}(\sigma)$ & $=\max_{j:\lambda_j>1}B_j$ [numerical threshold] & \S5.1 \\
\bottomrule
\end{tabular}
\end{center}

\section{Non-Circularity Remarks}

This paper uses no circular reasoning with respect to RH\@. The construction
of $T=\Phi^*\Phi$ from prime resonance vectors is explicit and postulates
neither GUE statistics, Montgomery's pair correlation, nor the Hilbert-P\'olya
conjecture. Geometric model content (\S\S\,7--10) is always marked as such
and is not used in any formal argument.

The following distinctions are maintained throughout:
\begin{itemize}
\item The identity $\Delta=\Estr-\Espec$ is an algebraic theorem
(Theorem~\ref{thm:identity}); the sign $\Delta\geq 0$ is numerical only ---
these are distinct levels.
\item The eigenvalues $\lambda_j$ of $T=\Phi^*\Phi$ are the squares of the
singular values of $\Phi$, not the zero ordinates $\gamma_k$ themselves.
\item No claim relating $\etaorig$ to RH is proved or used.
\end{itemize}

\section{Sample Output of \texttt{eta\_verification.py}}

The following output was produced with $\kappa=53$, $\varepsilon=0.05$,
$N=100$, $\sigma=0.5$:

\begin{verbatim}
-- NUM4 Benchmark (kappa=53, eps=0.05, sigma=0.5, N=100) --
eta_orig = 0.66926873
Reference:  0.66926893   (Delta = 1.97e-07)

-- NUM2 Normative table (kappa=53, eps=0.05, N=100) --------
 sigma | eta_orig
 0.10  | 0.70441977
 0.20  | 0.69826865
 0.30  | 0.69038398
 0.40  | 0.68072993
 0.50  | 0.66926873
 0.60  | 0.65602069
 0.70  | 0.64110102
 0.80  | 0.62472010
 0.90  | 0.60715692
 0.95  | 0.59802896

Monotone decreasing? True
\end{verbatim}

Requirements: Python~3.x, NumPy~$\geq 1.24$, mpmath~$\geq 1.3$.
Command: \texttt{python eta\_verification.py}

\bigskip
\noindent\textit{U. Tehrani \textperiodcentered\ March 2026}

\end{document}
